\input{sol_pre.tex}
\input{q_2026-10-03_1.tex}
\begin{document}
\input{sol_head.tex}

%================= 1
\Savol{1}
$36=4\cdot9$. \ 4 ga böliniş: oxirgi ikki raqam $\overline{5a}$ 4 ga bölinadi $\Rightarrow$ $a=2$ yoki $a=6$.\\
9 ga böliniş: $3+4+b+5+a=12+a+b$ 9 ga bölinadi.\\
$a=2$: $14+b\Rightarrow b=4$. \quad $a=6$: $18+b\Rightarrow b=0$ yoki $b=9$. \quad Jami 3 ta juftlik.
\Javob{C) 3 ta}

%================= 2
\Savol{2}
\[2^7+4^4+16^3=2^7+2^8+2^{12}=2^7(1+2+32)=2^7\cdot35=2^7\cdot5\cdot7.\]
\Javob{B) 7}

%================= 3
\Savol{3}
Tögarakka qatnaşmaydiganlar: matematikaga $100-84=16$, ingliz tiliga $100-68=32$, fizikaga $100-72=28$.\\
Kamida bittasiga qatnaşmaydiganlar köpi bilan $16+32+28=76$ ta. Demak uçalasiga qatnaşadiganlar kamida $100-76=24$ ta.
\Javob{B) 24}

%================= 4
\Savol{4}
Birinçi guruh yolğiz $x$ kunda, ikkinçisi $y$ kunda tugatsin. $13\frac13\%=\frac2{15}$:
\[0{,}4x-\tfrac2{15}y=2\ \Rightarrow\ y=3x-15;\qquad \frac1x+\frac1{3x-15}=\frac16.\]
\[6(3x-15)+6x=x(3x-15)\ \Rightarrow\ x^2-13x+30=0\ \Rightarrow\ x=10\ \ (x=3\ \text{da}\ y<0).\]
\Javob{C) 10}

%================= 5
\Savol{5}
Surat: $2^5\cdot3^8\cdot7^5+3^5\cdot(2^5\cdot7^5)\cdot(2^4\cdot3)=2^5\cdot3^6\cdot7^5(9+16)=25\cdot2^5\cdot3^6\cdot7^5$.\\
Maxraj: $5^4\cdot2\cdot7^2\cdot2^5\cdot3^5-2^5\cdot3^5\cdot5^2\cdot7^2=2^5\cdot3^5\cdot5^2\cdot7^2(50-1)=25\cdot2^5\cdot3^5\cdot7^4$.
\[\frac{25\cdot2^5\cdot3^6\cdot7^5}{25\cdot2^5\cdot3^5\cdot7^4}=3\cdot7=21.\]
\Javob{C) 21}

%================= 6
\Savol{6}
$n(n+1)(n+2)(n+3)+1=(n^2+3n)(n^2+3n+2)+1=(n^2+3n+1)^2$. \ $n=28$:
\[\sqrt{28\cdot29\cdot30\cdot31+1}=28^2+3\cdot28+1=784+84+1=869.\]
\Javob{C) 869}

%================= 7
\Savol{7}
$u=3\sqrt3+2\sqrt5$, \ $v=\sqrt{15}+6$: ifoda $(u+v)(u-v)=u^2-v^2$.
\[u^2=27+20+12\sqrt{15}=47+12\sqrt{15},\qquad v^2=15+36+12\sqrt{15}=51+12\sqrt{15},\qquad u^2-v^2=-4.\]
\Javob{D) $-4$}

%================= 8
\Savol{8}
$a_{11}$ — 21 ta hadning örtasi: \ $S_{21}=21\cdot a_{11}=21\cdot24=504$.
\Javob{A) 504}

%================= 9
\Savol{9}
$(x+7)^2=(x+3)(2x+8)\Rightarrow x^2+14x+49=2x^2+14x+24\Rightarrow x^2=25$.\\
$x=5$: $8,\ 12,\ 18$ — ösuvçi, $q=\frac32$. \quad $x=-5$: $-2,\ 2,\ -2$ — ösuvçi emas.
\[b_5=8\cdot\left(\tfrac32\right)^4=8\cdot\tfrac{81}{16}=40{,}5.\]
\Javob{D) 40,5}

%================= 10
\Savol{10}
$a^4-b^4=(a-b)(a+b)(a^2+b^2)$:
\[\frac{a(a^2-b^2)-b(a-b)^2}{(a-b)(a+b)(a^2+b^2)}=\frac{(a-b)\bigl(a^2+ab-ab+b^2\bigr)}{(a-b)(a+b)(a^2+b^2)}=\frac1{a+b}.\]
\Javob{B) $\frac1{a+b}$}

%================= 11
\Savol{11}
$x+6\pm6\sqrt{x-3}=(x-3)\pm6\sqrt{x-3}+9=\left(\sqrt{x-3}\pm3\right)^2$. \ $x>12$ da $\sqrt{x-3}>3$:
\[\left(\sqrt{x-3}+3\right)-\left(\sqrt{x-3}-3\right)=6.\]
\Javob{C) 6}

%================= 12
\Savol{12}
$4\sin2x-6\cos2x$ ning qiymatlar sohasi $\left[-\sqrt{16+36};\sqrt{16+36}\right]=[-\sqrt{52};\sqrt{52}]$, \ $\sqrt{52}\approx7{,}2$.\\
Butun $a$: $-7,\dots,7$ — 15 ta.
\Javob{B) 15 ta}

%================= 13
\Savol{13}
$\sin x\ne1$. \ $\sin x+1=\cos x(\sin x-1)\Rightarrow \sin x+\cos x-\sin x\cos x+1=0$.\\
$s=\sin x+\cos x$, \ $\sin x\cos x=\frac{s^2-1}2$: \ $2s-s^2+1+2=0\Rightarrow s^2-2s-3=0\Rightarrow s=-1$ ($s=3$ — mumkin emas).
\[\sqrt2\sin\left(x+\tfrac\pi4\right)=-1\ \Rightarrow\ x=-\tfrac\pi2+2\pi k\ \ \text{yoki}\ \ x=\pi+2\pi k.\]
$\left[\frac\pi2;\frac{5\pi}2\right]$ da: $\pi$ va $\frac{3\pi}2$; \ yiğindi $\frac{5\pi}2$.
\Javob{C) $\frac{5\pi}2$}

%================= 14
\Savol{14}
$25^x$ ga bölamiz, $t=\left(\frac35\right)^x>0$: \ $3t^2-7t+4=0\Rightarrow t=1$ yoki $t=\frac43$ — ikkalasi musbat, 2 ta ildiz ($x=0$, $x=\log_{0{,}6}\frac43$).
\Javob{D) 2 ta}

%================= 15
\Savol{15}
$x^2=|x|^2$: \ $|x|^2-3|x|+2=0\Rightarrow|x|=1$ yoki $|x|=2$ $\Rightarrow$ $x=\pm1,\ \pm2$ — 4 ta.
\Javob{A) 4 ta}

%================= 16
\Savol{16}
$|x-1|+|2-x|\ge|(x-1)+(2-x)|=1$, tenglik $(x-1)(2-x)\ge0$, ya'ni $x\in[1;2]$ da.\\
Demak yeçim $[1;2]$, butun sonlar $1;2$, yiğindi 3.
\Javob{B) 3}

%================= 17
\Savol{17}
$t=2025^x>0$: \ $t+\dfrac{2026\cdot2025}t=4051\Rightarrow t^2-4051t+2025\cdot2026=0\Rightarrow t=2025$ yoki $t=2026$ (Viyet).\\
$x=1$ va $x=\log_{2025}2026$ — 2 ta.
\Javob{A) 2 ta}

%================= 18
\Savol{18}
$2^x-1>0$, $2^x-1\ne1$ $\Rightarrow$ $x>0$, $x\ne1$. \ $u=\log_3(2^x-1)\ne0$: \ $u+\dfrac1u\le\dfrac{10}3$.\\
$u<0$ da doim bajariladi: $0<2^x-1<1\Rightarrow 0<x<1$ — butun son yöq.\\
$u>0$: $3u^2-10u+3\le0\Rightarrow\frac13\le u\le3\Rightarrow \sqrt[3]3\le2^x-1\le27\Rightarrow 1+\sqrt[3]3\le2^x\le28$.\\
$2^1=2<1+\sqrt[3]3\approx2{,}44$, \ $2^4=16\le28<2^5$: butun $x=2;3;4$, yiğindi $9$.
\Javob{A) 9}

%================= 19
\Savol{19}
$x\le1$ da öng tomon $\le0$ — tengsizlik bajariladi.\\
\textbf{$2\le x\le4$}: $|x^2-6x+8|=-x^2+6x-8\in[0;1]$, \ $\bigl||\dots|-3\bigr|=x^2-6x+11$: \ $x^2-7x+12\ge0\Rightarrow x\le3$ yoki $x\ge4$ $\Rightarrow$ $[2;3]\cup\{4\}$.\\
\textbf{$x<2$ yoki $x>4$}: $|x^2-6x+8|-3=(x-1)(x-5)$.\\
\quad $1<x<2$ yoki $4<x<5$: $(x-1)(5-x)\ge x-1\Rightarrow x\le4$ — $(1;2)$ töliq, $(4;5)$ dan heç narsa.\\
\quad $x\ge5$: $(x-1)(x-5)\ge x-1\Rightarrow x\ge6$.\\
Natija: $(-\infty;3]\cup\{4\}\cup[6;+\infty)$.
\Javob{B)}

%================= 20
\Savol{20}
Davr 3: \ $f(8)=f(2)$, \ $f(-4)=f(-4+6)=f(2)$, \ $f(-1)=f(2)$.
\[5+5-5=5.\]
\Javob{C) 5}

%================= 21
\Savol{21}
Grafik $(0;0)$ dan ötadi va $(2;0)$ da $Ox$ ga urinadi ($x=2$ — ikki karrali ildiz): \ $f(x)=ax(x-2)^2$.\\
Maksimum $\left(\frac23;\frac{32}{27}\right)$: \ $a\cdot\frac23\cdot\frac{16}9=\frac{32}{27}a=\frac{32}{27}\Rightarrow a=1$.
\[f(-2)=-2\cdot16=-32.\]
\Javob{A) $-32$}

%================= 22
\Savol{22}
$x+3=(x-4)+7$:
\[F(x)=\int\left((x-4)^{-\frac12}+7(x-4)^{-\frac32}\right)dx=2\sqrt{x-4}-\frac{14}{\sqrt{x-4}}+C.\]
$F(5)=2-14+C=0\Rightarrow C=12$. \quad $F(8)=4-7+12=9$.
\Javob{D) 9}

%================= 23
\Savol{23}
\fig{\begin{tikzpicture}[scale=0.85]
\coordinate (A) at (0,0); \coordinate (C) at (5,0); \coordinate (B) at (0.791,3.172); \coordinate (D) at (1.811,2.403);
\coordinate (F) at (0.672,2.696); \coordinate (E) at (0.672,0);
\coordinate (P) at (intersection of A--D and F--E);
\draw[thick] (A)--(B)--(C)--cycle; \draw[thick] (A)--(D); \draw[thick] (F)--(E);
\draw (0.672,0.2)--(0.872,0.2)--(0.872,0);
\pic[draw,angle radius=0.18cm]{right angle=A--D--C};
\pic[draw=aux,angle radius=0.55cm,"\scriptsize\color{aux}$53^\circ$",angle eccentricity=1.7] {angle=E--A--P};
\pic[draw=aux,angle radius=0.6cm,"\scriptsize\color{aux}$37^\circ$",angle eccentricity=1.6] {angle=B--C--A};
\pic[draw,angle radius=0.3cm,"\scriptsize$\alpha$",angle eccentricity=1.6] {angle=D--P--F};
\foreach \p in {A,B,C,D,E,F,P} \fill (\p) circle (1.2pt);
\node[below left] at (A){$A$}; \node[above] at (B){$B$}; \node[below right] at (C){$C$}; \node[above right] at (D){$D$};
\node[below] at (E){$E$}; \node[left] at (F){$F$}; \node[left] at (P){$P$};
\end{tikzpicture}}
$\angle C=180^\circ-76^\circ-67^\circ=37^\circ$. \ $\triangle ADC$ ($\angle D=90^\circ$): \ $\angle DAC=90^\circ-37^\circ=53^\circ$.\\
$\triangle APE$ ($\angle E=90^\circ$): \ $\angle APE=90^\circ-53^\circ=37^\circ$. \ $\alpha$ unga vertikal burçak:
\[\alpha=37^\circ.\]
\Javob{A) $37^\circ$}

%================= 24
\Savol{24}
$\varphi(x)=\dfrac{3+x}{1-x}$, \ $\psi(x)=\dfrac{x-3}{x+1}$. Tekşiriş: $\varphi(\varphi(x))=\dfrac{3+\frac{3+x}{1-x}}{1-\frac{3+x}{1-x}}=\dfrac{6-2x}{-2-2x}=\psi(x)$ va $\varphi(\psi(x))=x$.\\
Şartda $x$ ni $\varphi(x)$ va $\psi(x)$ bilan almaştiramiz:
\[f(\varphi)+f(\psi)=x,\qquad f(\psi)+f(x)=\varphi,\qquad f(x)+f(\varphi)=\psi.\]
Oxirgi ikkisini qöşib, birinçisini ayiramiz: \ $2f(x)=\varphi(x)+\psi(x)-x$.
\[f'(x)=\frac12\left(\frac4{(1-x)^2}+\frac4{(x+1)^2}-1\right),\qquad f'(0)=\frac12(4+4-1)=\frac72.\]
\Javob{B) $\frac72$}

%================= 25
\Savol{25}
Ştrixlangan soha $=S_{\triangle OAC}-S_{\text{sektor }OAB}$. \ $AC=OA\tg30^\circ=\dfrac{10}{\sqrt3}$:
\[S_{OAC}=\frac12\cdot10\cdot\frac{10}{\sqrt3}=\frac{50\sqrt3}3,\qquad S_{\text{sektor}}=\frac{30}{360}\pi\cdot100=\frac{25\pi}3.\]
\[S=\frac{50\sqrt3-25\pi}3.\]
\Javob{C) $\frac{50\sqrt3-25\pi}3$}

%================= 26
\Savol{26}
\fig{\begin{tikzpicture}[scale=0.55]
\coordinate (B) at (0,0); \coordinate (C) at (10,0); \coordinate (A) at (7.25,5.333); \coordinate (M) at (5,0); \coordinate (I) at (6,0); \coordinate (H) at (7.25,0);
\fill[msblue!25] (A)--(M)--(I)--cycle;
\draw[thick] (A)--(B)--(C)--cycle; \draw[thick] (A)--(M); \draw[thick] (A)--(I);
\draw[ax] (A)--(H); \draw[aux] (7.05,0)--(7.05,0.2)--(7.25,0.2);
\foreach \p in {A,B,C,M,I,H} \fill (\p) circle (2pt);
\node[above] at (A){$A$}; \node[below left] at (B){$B$}; \node[below right] at (C){$C$}; \node[below] at (M){$M$}; \node[below] at (I){$I$}; \node[aux,below] at (H){$H$};
\node[above left] at (3.6,2.7){\small$9$}; \node[above right] at (8.6,2.7){\small$6$}; \node[below] at (2.5,0){\small$5$}; \node[aux,right] at (7.25,2.6){\small$h$};
\end{tikzpicture}}
$AB=9$, $AC=6$, $BM=5\Rightarrow BC=10$. Bissektrisa xossasi: $\dfrac{BI}{IC}=\dfrac96\Rightarrow BI=6$, \ $MI=1$.\\
$\triangle AMI$ va $\triangle ABC$ ning balandligi $AH$ umumiy (\textcolor{aux}{qizil}), şuning uçun $S_{AMI}=\dfrac{MI}{BC}S_{ABC}=\dfrac1{10}S_{ABC}$.\\
Geron formulasi ($p=12{,}5$):
\[S_{ABC}=\sqrt{12{,}5\cdot2{,}5\cdot3{,}5\cdot6{,}5}=\sqrt{\frac{25\cdot5\cdot7\cdot13}{16}}=\frac{5\sqrt{455}}4,\qquad S_{AMI}=\frac{\sqrt{455}}8.\]
\Javob{A) $\frac{\sqrt{455}}8$}

%================= 27
\Savol{27}
\fig{\begin{tikzpicture}[scale=0.42]
\coordinate (A) at (-6,0); \coordinate (D) at (6,0); \coordinate (B) at (-2,8); \coordinate (C) at (2,8); \coordinate (O) at (0,6);
\coordinate (P) at (-4,4); \coordinate (E) at (-2,4); \coordinate (F) at (2,4);
\fill[msblue!25] (A)--(P)--(O)--cycle;
\draw[thick] (A)--(B)--(C)--(D)--cycle; \draw[thick] (A)--(C); \draw[thick] (B)--(D); \draw[thick] (P)--(O); \draw[thick] (E)--(F);
\draw[ax] (0,8)--(0,0); \draw[ax] (C)--(2,0);
\pic[draw,angle radius=0.25cm]{right angle=C--O--D};
\foreach \p in {A,B,C,D,O,P,E,F} \fill (\p) circle (4pt);
\node[below left] at (A){$A$}; \node[below right] at (D){$D$}; \node[above left] at (B){$B$}; \node[above right] at (C){$C$};
\node[above] at (0,6.25){$O$}; \node[left] at (P){$P$}; \node[below] at (E){$E$}; \node[below] at (F){$F$};
\node[aux,right] at (2,2){\small$h$};
\end{tikzpicture}}
$EF$ — trapetsiya diagonallari örtalarini tutaştiruvçi kesma: $EF=\dfrac{AD-BC}2=4\Rightarrow AD-BC=8$.\\
Teng yonli trapetsiyada diagonallar perpendikulyar bölsa, $\triangle AOD$, $\triangle BOC$ teng yonli töğri burçakli va balandlik $h=\dfrac{AD+BC}2$ (\textcolor{aux}{qizil}).\\
Yon tomon: $CD^2=h^2+\left(\frac{AD-BC}2\right)^2\Rightarrow80=h^2+16\Rightarrow h=8$, \ $AD+BC=16$ $\Rightarrow$ $AD=12$, $BC=4$.\\
$OA=\dfrac{AD}{\sqrt2}=6\sqrt2$, \ $OB=\dfrac{BC}{\sqrt2}=2\sqrt2$, \ $S_{AOB}=\frac12\cdot6\sqrt2\cdot2\sqrt2=12$. $OP$ — mediana:
\[S_{APO}=\frac12S_{AOB}=6.\]
\Javob{D) 6}

%================= 28
\Savol{28}
\fig{\begin{tikzpicture}[scale=0.9]
\coordinate (A) at (0,0); \coordinate (E) at (2,0); \coordinate (D) at (2.618,1.902); \coordinate (C) at (1,3.078); \coordinate (B) at (-0.618,1.902);
\coordinate (R) at (3.236,0); \coordinate (T) at (4.236,3.078);
\fill[aux!10] (A)--(C)--(R)--cycle;
\draw[thick] (A)--(B)--(C)--(D)--(E)--cycle; \draw[thick] (A)--(C)--(T)--(R)--cycle; \draw[thick] (C)--(R);
\pic[draw=aux,angle radius=0.6cm,"\scriptsize\color{aux}$72^\circ$",angle eccentricity=1.45] {angle=R--A--C};
\pic[draw=aux,angle radius=0.75cm,"\scriptsize\color{aux}$54^\circ$",angle eccentricity=1.4] {angle=A--C--R};
\pic[draw,angle radius=1.2cm,"\scriptsize$x$",angle eccentricity=1.2] {angle=R--C--D};
\foreach \p in {A,B,C,D,E,R,T} \fill (\p) circle (1.2pt);
\node[below left] at (A){$A$}; \node[left] at (B){$B$}; \node[above] at (C){$C$}; \node[right] at (D){$D$}; \node[below] at (E){$E$}; \node[below] at (R){$R$}; \node[above right] at (T){$T$};
\end{tikzpicture}}
Muntazam beşburçak burçagi $108^\circ$. $\triangle ABC$ teng yonli: $\angle BAC=\angle BCA=36^\circ$, \ $\angle CAE=108^\circ-36^\circ=72^\circ$, \ $\angle ACD=108^\circ-36^\circ=72^\circ$.\\
Romb tomoni $AR=AC$ $\Rightarrow$ $\triangle ACR$ teng yonli (\textcolor{aux}{qizil}): \ $\angle ACR=\dfrac{180^\circ-72^\circ}2=54^\circ$.
\[x=\angle ACD-\angle ACR=72^\circ-54^\circ=18^\circ.\]
\Javob{A) $18^\circ$}

%================= 29
\Savol{29}
\fig{\begin{tikzpicture}[scale=0.4]
\draw[thick] (-1,0)--(9,0) node[right]{\small$\alpha$};
\coordinate (A1) at (0,0); \coordinate (B1) at (7,0); \coordinate (C1) at (3,0); \coordinate (A) at (0,2.25); \coordinate (B) at (7,4.75); \coordinate (C) at (3,3.321);
\draw[thick] (A)--(B); \draw[thick] (A)--(A1); \draw[thick] (B)--(B1); \draw[axs] (C)--(C1);
\foreach \p in {A,B,C,A1,B1,C1} \fill (\p) circle (4pt);
\node[left] at (A){$A$}; \node[right] at (B){$B$}; \node[above] at (C){$C$}; \node[below] at (A1){$A_1$}; \node[below] at (B1){$B_1$}; \node[aux,below] at (C1){$C_1$};
\node[left] at (0,1.1){\small$9$}; \node[right] at (7,2.4){\small$19$};
\end{tikzpicture}}
$AA_1\parallel CC_1\parallel BB_1$ (hammasi tekislikka perpendikulyar), $AA_1B_1B$ — trapetsiya. $C$ asoslarga parallel $CC_1$ kesmani $AC:CB=3:4$ nisbatda böladi:
\[CC_1=\frac{4\cdot AA_1+3\cdot BB_1}{7}=\frac{36+57}7=\frac{93}7\ \text{sm}.\]
\Javob{A) $\frac{93}7$ sm}

%================= 30
\Savol{30}
$\vec a\cdot\vec b=2\cdot4\cdot\cos60^\circ=4$, \ $\vec a^{\,2}=4$, \ $\vec b^{\,2}=16$.
\[(2\vec a-3\vec b)(2\vec a-5\vec b)=4\vec a^{\,2}-16\,\vec a\cdot\vec b+15\vec b^{\,2}=16-64+240=192.\]
\Javob{B) 192}

%================= 31
\Savol{31}
$x^2>2025\iff|x|>45$; \ $|x|\le50$. \ $A\cap B=\{\pm46,\pm47,\pm48,\pm49,\pm50\}$ — 10 ta element.
\[2^{10}=1024.\]
\Javob{C) 1024}

%================= 32
\Savol{32}
Koeffitsiyentlar yiğindisi — köphadning $x=y=1$ dagi qiymati: $(1+2)^5=243$.
\Javob{D) 243}

%================= 33-35
\newpage
\Savol{33}
\fig{\begin{tikzpicture}[scale=0.6]
\draw[thick] (0,0) circle (1.6); \draw[thick] (-1.131,-1.131) rectangle (1.131,1.131);
\draw[axs] (0,0)--(1.131,1.131); \draw[ax] (0,0)--(1.131,0);
\fill (0,0) circle (1.3pt);
\node[aux] at (0.35,0.65){\scriptsize$R_1$}; \node[aux,below] at (0.6,0){\scriptsize$r$};
\node[below] at (0,-1.7){\small Silindr};
\begin{scope}[xshift=4.5cm]
\draw[thick] (0,0) circle (1.6); \draw[thick] (-1.386,-0.8)--(0,1.6)--(1.386,-0.8)--cycle;
\draw[axs] (0,0)--(1.386,-0.8); \draw[ax] (0,1.6)--(0,-0.8);
\fill (0,0) circle (1.3pt);
\node[aux] at (0.75,-0.15){\scriptsize$R_2$}; \node[aux,below] at (0.7,-0.8){\scriptsize$\rho$};
\node[below] at (0,-1.7){\small Konus};
\end{scope}
\end{tikzpicture}}
\textbf{Silindr} (radiusi $r$, balandligi $2r$): öq kesimi — kvadrat, uning diagonali şar diametri: $2r\sqrt2=2R_1\Rightarrow R_1=r\sqrt2$.\\
$S_{\text{sil}}=2\pi r^2+2\pi r\cdot2r=6\pi r^2$, \ $V_{\text{sil}}=2\pi r^3$.\\
\textbf{Konus} (radiusi $\rho$, yasovçisi $2\rho$, balandligi $\rho\sqrt3$): öq kesimi — tomoni $2\rho$ teng tomonli uçburçak, $R_2=\dfrac{2\rho}{\sqrt3}$.\\
$S_{\text{kon}}=\pi\rho^2+\pi\rho\cdot2\rho=3\pi\rho^2$, \ $V_{\text{kon}}=\frac13\pi\rho^2\cdot\rho\sqrt3=\frac{\sqrt3}3\pi\rho^3$.

\textbf{33.} $6\pi r^2=3\pi\rho^2\Rightarrow\dfrac r\rho=\dfrac1{\sqrt2}=\dfrac{\sqrt2}2$.

\textbf{34.} $\dfrac{V_{\text{sil}}}{V_{\text{kon}}}=\dfrac{2r^3}{\frac{\sqrt3}3\rho^3}=\dfrac6{\sqrt3}\cdot\dfrac1{2\sqrt2}=\dfrac{\sqrt3}{\sqrt2}=\dfrac{\sqrt6}2$.

\textbf{35.} $\dfrac{R_1}{R_2}=\dfrac{r\sqrt2}{2\rho/\sqrt3}=\dfrac{\sqrt6}2\cdot\dfrac r\rho=\dfrac{\sqrt6}2\cdot\dfrac1{\sqrt2}=\dfrac{\sqrt3}2$.
\Javob{33 – B; \ 34 – C; \ 35 – A}

%================= 36
\newpage
\Savol{36}
Ikkinçi tenglama: $2^{\frac y2}=2^{2|x|+2}\Rightarrow y=4|x|+4$.\\
\textbf{Zaruriy şart.} Agar $(x;y)$ yeçim bölsa, $(-x;y)$ ham yeçim (sistemada $x$ faqat $x^2$, $|x|$ körinişida). Yagona bölişi uçun $x=0$, unda $y=4$:
\[\log_3(3a^2)=1+\log_31\ \Rightarrow\ 3a^2=3\ \Rightarrow\ a=\pm1.\]
\textbf{Yetarlilik.} $a=1$: \ $x^4y^2+3=3-3x^2y^2\Rightarrow x^2y^2(x^2+3)=0$; \ $y\ge4$ bölgani uçun faqat $x=0$ — yagona yeçim $(0;4)$.\\
$a=-1$: \ $x^4y^2+3=3+3x^2y^2\Rightarrow x^2y^2(x^2-3)=0\Rightarrow x=0$ yoki $x=\pm\sqrt3$ — 3 ta yeçim, mos emas.

a) Yagona qiymat $a=1$.

b) $x+y=0+4=4$.
\Javob{a) 1; \quad b) 4}

%================= 37
\newpage
\Savol{37}
Ikkala tomonni 2 ga bölamiz:
\[\tfrac12\sin x+\tfrac{\sqrt3}2\cos x=\tfrac12\ \Rightarrow\ \sin\left(x+\tfrac\pi3\right)=\tfrac12.\]
\[x+\tfrac\pi3=\tfrac\pi6+2\pi k\Rightarrow x=-\tfrac\pi6+2\pi k;\qquad x+\tfrac\pi3=\tfrac{5\pi}6+2\pi k\Rightarrow x=\tfrac\pi2+2\pi k.\]
a) $[-2\pi;2\pi]$ da: $-\frac\pi6,\ \frac{11\pi}6,\ \frac\pi2,\ -\frac{3\pi}2$.
\[-\frac\pi6+\frac{11\pi}6+\frac\pi2-\frac{3\pi}2=\frac{5\pi}3-\pi=\frac{2\pi}3.\]
b) $[0;5\pi]$ da: $\frac{11\pi}6,\ \frac{23\pi}6$ \ va \ $\frac\pi2,\ \frac{5\pi}2,\ \frac{9\pi}2$ — 5 ta.
\Javob{a) $\frac{2\pi}3$; \quad b) 5}

%================= 38
\newpage
\Savol{38}
\fig{\begin{tikzpicture}[xscale=0.45,yscale=0.22]
\draw[->] (-0.5,0)--(12,0) node[right]{$x$}; \draw[->] (0,-2)--(0,11) node[above]{$y$};
\draw[thick,domain=1:11,samples=60] plot (\x,{0.2887*(\x-6)*(\x-6)+2.598});
\draw[thick,domain=1.2:7] plot (\x,{-1.732*(\x-6)}); \draw[thick,domain=5:10.8] plot (\x,{1.732*(\x-6)});
\draw[ax] (6,0)--(6,2.598); \draw[ax] (3,0)--(3,5.196); \draw[ax] (9,0)--(9,5.196);
\fill (3,5.196) circle (3pt); \fill (9,5.196) circle (3pt); \fill[aux] (6,2.598) circle (3pt);
\node[below] at (3,0){$3$}; \node[below] at (6.3,0){$6$}; \node[below] at (9,0){$9$};
\node[aux] at (5.0,0.7){\scriptsize$60^\circ$}; \node[aux] at (7.0,0.7){\scriptsize$60^\circ$};
\node[aux,right] at (6.1,3.4){\small$\left(6;\frac{3\sqrt3}2\right)$};
\end{tikzpicture}}
Yoyiq burçak $3\beta=180^\circ\Rightarrow\beta=60^\circ$. Urinmalar $Ox$ bilan $60^\circ$ li burçak hosil qiladi:
\[k_1=-\tg60^\circ=-\sqrt3,\qquad k_2=\sqrt3,\qquad y_1=-\sqrt3(x-6),\quad y_2=\sqrt3(x-6).\]
Uriniş şartlari: $y'(3)=6a+b=-\sqrt3$, \ $y'(9)=18a+b=\sqrt3$ $\Rightarrow$ $a=\dfrac{\sqrt3}6$, \ $b=-2\sqrt3$.

a) $\dfrac{ab}{k_1k_2}=\dfrac{\frac{\sqrt3}6\cdot(-2\sqrt3)}{-3}=\dfrac{-1}{-3}=\dfrac13$.
b) Uç abssissasi $x_0=-\dfrac b{2a}=6$. \ $y=a(x-6)^2+m$, \ $y(3)=y_1(3)=3\sqrt3$:
\[\frac{\sqrt3}6\cdot9+m=3\sqrt3\ \Rightarrow\ m=3\sqrt3-\frac{3\sqrt3}2=\frac{3\sqrt3}2.\]
\Javob{a) $\frac13$; \quad b) $\frac{3\sqrt3}2$}

%================= 39
\newpage
\Savol{39}
a) Urinma $(-4;0)$ va $(0;2)$ dan ötadi: burçak koeffitsiyenti $f'(2)=\dfrac{2-0}{0-(-4)}=\dfrac12$.\\
Urinma tenglamasi $y=\frac12x+2$, uriniş nuqtasida $f(2)=\frac12\cdot2+2=3$.

b) Bölinma hosilasi:
\[h'(x)=\frac{f'(x)(x+2)-f(x)}{(x+2)^2},\qquad h'(2)=\frac{\frac12\cdot4-3}{16}=-\frac1{16}.\]
\Javob{a) $\frac12$; \quad b) $-\frac1{16}$}

%================= 40
\newpage
\Savol{40}
\fig{\begin{tikzpicture}[xscale=0.8,yscale=0.33]
\fill[msblue!20] plot[domain=-1:0,samples=30] (\x,{\x*\x*\x-3*\x*\x+5})--(3,5)--plot[domain=3:2,samples=30] (\x,{\x*\x*\x-3*\x*\x+5})--cycle;
\draw[->] (-2,0)--(4,0) node[right]{$x$}; \draw[->] (0,-1.2)--(0,7.5) node[above]{$y$};
\draw[thick,domain=-1.25:3.3,samples=80] plot (\x,{\x*\x*\x-3*\x*\x+5});
\draw[ax] (-1.8,1)--(3.6,1) node[right]{\small$y=1$}; \draw[ax] (-1.8,5)--(3.6,5) node[right]{\small$y=5$};
\foreach \p in {(-1,1),(2,1),(0,5),(3,5)} \fill \p circle (2pt);
\node[below left] at (-1,1){\small$-1$}; \node[below] at (2,0.95){\small$2$}; \node[above right] at (3,5){\small$3$};
\end{tikzpicture}}
$(-1;1)$ da kesişadi va $(2;1)$ da urinadi: $y-1$ ning ildizlari $-1$ (oddiy) va $2$ (ikki karrali):
\[y-1=(x+1)(x-2)^2=x^3-3x^2+4\ \Rightarrow\ y=x^3-3x^2+5,\quad a=-3,\ b=0,\ c=5.\]
$y'=3x^2-6x=0\Rightarrow x=0$ (maksimum), $k=y(0)=5$.

a) $\dfrac{a+b+c}k=\dfrac{-3+0+5}5=\dfrac25$.
b) $y=5$ çiziq grafikni yana $x^3-3x^2=0\Rightarrow x=3$ da kesadi. Soha egri çiziq bilan ikki qismga bölinadi:
\[S_1=\int_{-1}^{2}\bigl(x^3-3x^2+5-1\bigr)dx=\left(\frac{x^4}4-x^3+4x\right)\Big|_{-1}^{2}=4-\left(-\frac{11}4\right)=\frac{27}4,\]
\[S_2=\int_0^3\bigl(5-(x^3-3x^2+5)\bigr)dx=\left(x^3-\frac{x^4}4\right)\Big|_0^3=27-\frac{81}4=\frac{27}4.\]
\[S=\frac{27}4+\frac{27}4=\frac{27}2.\]
\Javob{a) $\frac25$; \quad b) $\frac{27}2$}

%================= 41
\newpage
\Savol{41}
\fig{\begin{tikzpicture}[scale=0.62]
\coordinate (B) at (0,0); \coordinate (C) at (8,0); \coordinate (A) at (2.5,4.330); \coordinate (I) at (3,1.732); \coordinate (O) at (4,0.577);
\draw[thick] (A)--(B)--(C)--cycle; \draw[thick] (I) circle (1.732); \draw[thick] (A)--(I)--(C);
\draw[axs] (B)--(I); \draw[axs] (I)--(O);
\pic[draw,angle radius=0.35cm,"\scriptsize$120^\circ$",angle eccentricity=2.0] {angle=C--I--A};
\pic[draw=aux,angle radius=0.6cm,"\scriptsize\color{aux}$60^\circ$",angle eccentricity=1.5] {angle=C--B--A};
\foreach \p in {A,B,C,I,O} \fill (\p) circle (2pt);
\node[above] at (A){$A$}; \node[below left] at (B){$B$}; \node[below right] at (C){$C$}; \node[left] at (2.85,1.95){$O_1$}; \node[aux,below right] at (O){$O_2$};
\node[above right] at (5.3,2.2){\small$7$}; \node[below] at (4,-0.05){\small$8$}; \node[left] at (1.25,2.3){\small$c$};
\end{tikzpicture}}
$O_1$ — bissektrisalar kesişiş nuqtasi: $\angle AO_1C=180^\circ-\dfrac{\angle A+\angle C}2=90^\circ+\dfrac{\angle B}2=120^\circ\Rightarrow\angle B=60^\circ$.\\
$AB=c$, kosinuslar teoremasi: $7^2=c^2+8^2-2\cdot8c\cdot\frac12\Rightarrow c^2-8c+15=0\Rightarrow c=5$ yoki $c=3$.\\
$c=3$ da $\cos A=\dfrac{9+49-64}{42}<0$ — ötmas uçburçak, demak $c=5$.

a) \[S=\frac12\cdot5\cdot8\cdot\sin60^\circ=10\sqrt3.\]
b) $R=\dfrac{AC}{2\sin B}=\dfrac7{\sqrt3}$, \ $p=\dfrac{5+7+8}2=10$, \ $r=\dfrac Sp=\sqrt3$. Eyler formulasi:
\[O_1O_2^2=R^2-2Rr=\frac{49}3-2\cdot\frac7{\sqrt3}\cdot\sqrt3=\frac{49}3-14=\frac73\ \Rightarrow\ O_1O_2=\frac{\sqrt{21}}3.\]
\Javob{a) $10\sqrt3$; \quad b) $\frac{\sqrt{21}}3$}

%================= 42
\newpage
\Savol{42}
\fig{\begin{tikzpicture}[scale=0.24]
\coordinate (A) at (-6.5,0); \coordinate (B) at (6.5,0); \coordinate (C) at (1.5,12); \coordinate (D) at (-1.5,12);
\coordinate (E) at (-2.5,6); \coordinate (F) at (2.5,6); \coordinate (Q) at (0,4.333); \coordinate (T) at (0,9.75);
\fill[msblue!25] (E)--(Q)--(F)--(T)--cycle;
\draw[thick] (A)--(B)--(C)--(D)--cycle; \draw[thick] (A)--(F); \draw[thick] (B)--(E); \draw[thick] (D)--(F); \draw[thick] (C)--(E);
\draw[ax] (-7,6)--(7,6); \draw[ax] (0,0)--(0,12); \draw[axs] (F)--(2.5,0);
\pic[draw,angle radius=0.25cm]{right angle=A--F--D};
\foreach \p in {E,F,Q,T} \fill (\p) circle (4pt);
\node[below left] at (A){$A$}; \node[below right] at (B){$B$}; \node[above right] at (C){$C$}; \node[above left] at (D){$D$};
\node[left] at (-2.6,6.4){$E$}; \node[right] at (2.6,6.4){$F$}; \node[below right] at (Q){$Q$}; \node[above right] at (T){$T$};
\node[aux,right] at (2.5,3){\small$6$};
\node[left] at (-4.2,6.5){\small$13$}; \node[above] at (0,12){\small$3$}; \node[below] at (-3,0){\small$13$};
\end{tikzpicture}}
Teng yonli trapetsiya: $AB=13$, $CD=3$, yon tomonlar 13. $D$ dan $AB$ ga tuşirilgan balandlik asosi $A$ dan $\frac{13-3}2=5$ uzoqlikda:
$h=\sqrt{13^2-5^2}=12$, \[\cos A=\frac5{13},\quad \cos\frac A2=\frac3{\sqrt{13}},\quad\sin\frac A2=\frac2{\sqrt{13}},\quad\tg\frac A2=\frac23.\]
$\angle A+\angle D=180^\circ$ $\Rightarrow$ bissektrisalar $F$ da perpendikulyar kesişadi: $AF=AD\cos\frac A2=3\sqrt{13}$.\\
$F$ dan $AB$ gaça masofa $AF\sin\frac A2=6$, $AB$ böylab $A$ dan masofa $AF\cos\frac A2=9$ (\textcolor{aux}{qizil}).

a) $A$ dan simmetriya öqigaça $6{,}5$, $F$ esa $9$ da: öqdan $2{,}5$ öngda. Simmetriya böyiça $E$ öqdan $2{,}5$ çapda, ikkalasi bir balandlikda:
$EF=2{,}5+2{,}5=5$.
b) Böyalgan $EQFT$ — simmetriya öqiga nisbatan simmetrik deltoid, diagonallari $EF=5$ va $QT$.\\
$Q$ — $AF$ va $BE$ ning öqdagi kesişişi: \ $Q$ balandligi $6{,}5\cdot\tg\frac A2=6{,}5\cdot\frac23=\frac{13}3$.\\
$T$ — $DF$ va $CE$ ning kesişişi. $D$ dagi bissektrisa $DC$ bilan $\frac{\angle D}2=90^\circ-\frac A2$ burçak hosil qiladi, uning qiyaligi $\ctg\frac A2=\frac32$. $D$ dan öqqaça $1{,}5$ birlik: pasayiş $1{,}5\cdot\frac32=2{,}25$, $T$ balandligi $12-2{,}25=9{,}75$.
$QT=9{,}75-\dfrac{13}3=\dfrac{65}{12}$, \quad $S=\dfrac12\cdot5\cdot\dfrac{65}{12}=\dfrac{325}{24}$.
\Javob{a) 5; \quad b) $\frac{325}{24}$}

%================= 43
\newpage
\Savol{43}
\fig{\begin{tikzpicture}[scale=0.24]
\coordinate (A) at (0,0); \coordinate (B) at (6,8); \coordinate (C) at (20,3); \coordinate (D) at (22,-11); \coordinate (E) at (6,-10.392);
\coordinate (M) at ($(A)!0.5!(B)$); \coordinate (K) at ($(B)!0.5!(C)$); \coordinate (N) at ($(C)!0.5!(D)$); \coordinate (L) at ($(D)!0.5!(E)$);
\coordinate (P) at ($(M)!0.5!(N)$); \coordinate (T) at ($(K)!0.5!(L)$); \coordinate (H) at ($(A)!(P)!(E)$);
\fill[msblue!25] (A)--(P)--(T)--(E)--cycle;
\draw[thick] (A)--(B)--(C)--(D)--(E)--cycle; \draw[thick] (M)--(N); \draw[thick] (K)--(L); \draw[thick] (A)--(P)--(T)--(E);
\draw[axs] (P)--(H); \pic[draw=aux,angle radius=0.2cm]{right angle=P--H--A};
\pic[draw,angle radius=0.6cm,"\scriptsize$60^\circ$",angle eccentricity=1.6] {angle=E--A--P};
\foreach \p in {A,B,C,D,E,M,K,N,L,P,T} \fill (\p) circle (5pt);
\node[left] at (A){$A$}; \node[above] at (B){$B$}; \node[right] at (C){$C$}; \node[right] at (D){$D$}; \node[below] at (E){$E$};
\node[above left] at (M){$M$}; \node[above] at (K){$K$}; \node[right] at (N){$N$}; \node[below] at (L){$L$}; \node[above] at (P){$P$}; \node[right] at (T){$T$};
\node[aux] at (4.4,-2.5){\scriptsize$H$};
\node[above] at (6,0){\small$12$}; \node[left] at (2.6,-5.6){\small$12$};
\end{tikzpicture}}
Vektorlar bilan: $\vec P=\dfrac{\vec M+\vec N}2=\dfrac{\vec A+\vec B+\vec C+\vec D}4$, \ $\vec T=\dfrac{\vec K+\vec L}2=\dfrac{\vec B+\vec C+\vec D+\vec E}4$.
\[\overrightarrow{PT}=\vec T-\vec P=\frac{\vec E-\vec A}4=\frac14\overrightarrow{AE}.\]
a) Demak $PT\parallel AE$ va $PT=\dfrac{AE}4=3$.

b) $APTE$ — trapetsiya ($PT\parallel AE$), balandligi $PH=AP\sin60^\circ=6\sqrt3$ (\textcolor{aux}{qizil}):
\[S=\frac{AE+PT}2\cdot PH=\frac{12+3}2\cdot6\sqrt3=45\sqrt3.\]
\Javob{a) 3; \quad b) $45\sqrt3$}

%================= 44
\newpage
\Savol{44}
\fig{\begin{tikzpicture}[scale=0.8]
\draw[thick] (0,0) circle (2); \draw[thick] (-1.414,-1.414) rectangle (1.414,1.414);
\draw[axs] (0,0)--(1.414,1.414); \draw[ax] (0,0)--(1.414,0); \draw[ax] (0,0)--(0,1.414);
\fill (0,0) circle (1.5pt);
\node[aux] at (0.55,0.95){\scriptsize$R$}; \node[aux,below] at (0.7,0){\scriptsize$r$}; \node[aux,left] at (0,0.7){\scriptsize$\frac h2$};
\end{tikzpicture}}
Öq kesimi: silindr radiusi $r$, balandligi $h$; \ $r^2+\left(\dfrac h2\right)^2=R^2=4$ (\textcolor{aux}{qizil} töğri burçakli uçburçak).\\
Yon sirt $S=2\pi rh=4\pi\cdot r\cdot\dfrac h2$. Örta arifmetik va örta geometrik haqidagi tengsizlik:
\[r\cdot\frac h2\le\frac{r^2+\frac{h^2}4}2=2,\qquad\text{tenglik } r=\frac h2=\sqrt2\ \text{da}.\]
a) $S_{\max}=4\pi\cdot2=8\pi$.

b) $h=2\sqrt2$: \ $\dfrac Rh=\dfrac2{2\sqrt2}=\dfrac{\sqrt2}2$.
\Javob{a) $8\pi$; \quad b) $\frac{\sqrt2}2$}

%================= 45
\newpage
\Savol{45}
Silindr balandligi: $2\pi\cdot4\cdot H=80\pi\Rightarrow H=10$.\\
$V_{\text{sil}}=\pi\cdot16\cdot10=160\pi$, \quad $V_{\text{şar}}=\frac43\pi\cdot27=36\pi$.

a) \[\frac{V_{\text{şar}}}{V_{\text{sil}}}=\frac{36\pi}{160\pi}=\frac9{40}.\]
b) Şar botirilganda uning hajmiça suv tökildi. Şar olingaç, şu hajmdagi böşliq qoladi:
\[\pi\cdot4^2\cdot x=36\pi\ \Rightarrow\ x=\frac{36}{16}=\frac94.\]
\Javob{a) $\frac9{40}$; \quad b) $\frac94$}

\ShowJavoblar
\end{document}
